\documentclass[10pt]{article}

\usepackage[utf8]{inputenc}

\usepackage[T1]{fontenc}

\usepackage{amsmath}

\usepackage{amsfonts}

\usepackage{amssymb}

\usepackage[version=4]{mhchem}

\usepackage{stmaryrd}

\usepackage{bbold}

\usepackage{graphicx}

\usepackage[export]{adjustbox}

\graphicspath{ {./images/} }



\begin{document}

Ф. т. справедлива, в частности, для случая, когда $\mu_{x}$, $\mu_{y}$ и $\mu$-меры Лебега соответственно в евклидовых пространствах $\mathbb{R}^{m}, \mathbb{R}^{n}$ и $\mathbb{R} \boldsymbol{m}+\boldsymbol{n}$ ( $m$ и $n$-натуральные числа), $X=\mathbb{R}^{m}, Y=\mathbb{R}^{n}, X \times Y=\mathbb{R}^{m} \times \mathbb{R}^{n}=\mathbb{R}^{m+n}, f=f(x, y)$ измеримая по Лебегу на пространстве $\mathbb{R} m+n$ функция, $x \in \mathbb{R}^{m}, y \in \mathbb{R}^{n}$. При этих предшоложениях формула (1) имеет вид



\[

\iint_{\mathbb{R}^{m+n}} f(x, y) d x d y=\int_{\mathbb{R}^{n}} d y \int_{\mathbb{R}^{m}} f(x, y) d x .

\]



Для того чтобы в случае функции $f$, определенной на произвольном измеримом по Лебегу множестве $E \subset \mathbb{R}^{m+n}$, выразить кратный интеграл через повторный, нужно продолжить функцию $f$ нулем на все пространство $\mathbb{R}^{m+n}$ и применить формулу (2). См. также Повторный интеарал.



Ф. т. установлена Г. Фубини [1]



Jum.: [1] F u b i n i G., Sugli integrali multipli (1907), Opere scelte, v 2, Roma, 1958, p. 243-49. Л. Д. Кудрлвчев.



ФУБИНЙ ФОРМА - дифференциальная форма (квадратичная $F_{2}$ и кубическая $F_{3}$ ), на основе к-рой строится проективная дифференциальная геометрия. Введены Г. Фубпни (см. [1]).



\includegraphics[max width=\textwidth, center]{2023_07_09_8e5e0b3fc949bc32d289g-1}

динаты точки поверхности с внутренними координатами $u^{1}, u^{2}$ и пусть



\[

\begin{gathered}

b_{i j}=\left(x^{\alpha}, x_{1}^{\alpha}, x_{2}^{\alpha}, x_{i j}^{\alpha}\right), \\

b_{i j k}=a_{i j k}-\frac{1}{2} \partial_{k} b_{i j}+\frac{1}{2} b_{i, j} \partial_{k} \ln \sqrt{|b|}, \\

b=\operatorname{det} b_{i j}, \quad a_{i j k}=\left(x^{\alpha}, x_{1}^{\alpha}, x_{2}^{\alpha}, x_{i j k}^{\alpha}\right) .

\end{gathered}

\]



Тогда Ф. ф. определяются так:



\[

F_{2}=b_{i j} d u^{i} d u^{j}|b|^{-1 / 4}, \quad F_{3}=b_{i j k} d u^{i} d u^{j} d u^{k}|b|^{-1 / 4} \text {. }

\]



Однако сами проективные координаты не вполне определены: они допускают введение произвольных множителей и однородных линейных преобразований. Поәтому Ф. Ф. определены только с точностью до множителя и чтобы избежать связанных с этим затруднений, нормируют координаты и определенные через них формы; напр., при унимодулярных проективных преобразованиях Ф. Ф. сохраняют свое значение (с точностью до знака). Отношение $F_{3} / F_{2}$, наз. п р о е к т и в ны м лин е й ны м ә лем е н т о м, уже не зависит от нормирования (и определяет проективный метрич. әлемент).



Построенные метрич. средствами, исходя из второй квадратичной формы и формы Дарбу (определяемой Дарбу тензором), Ф. ф.



\[

f_{2}=\lambda F_{2}, \quad f_{3}=\lambda F_{3},

\]



инварианты относительно әквиаффинных преобразований и потому могут быть положены в основу эквиаффинной диффференциальной геометрии.



Jum.: [1] F u b i n i G., C e c h E., Geometria proiettiva differenzial, v. 1-2, Bologna, 1926-27; [2] K a $\mathrm{T}$ a $\mathrm{H}$ B. $\Phi$., Основы теории поверхностей..., т. 2, М.-Л., 1948; [3] ШШ и р ок о в П. А., ШІ и р о к о в А. П., Аффинная дифференциаль-

ная геометрия, М., 1959 .

М. В. Воицеховский. ФУБИНИ - IїТУДИ МЕТРИКА - эрмитова метрика на комплексном проективном пространстве $\mathbb{C} P^{n}$, определяемая әрмитовым скалярным произведением $(u, v)$ в пространстве $\mathbb{C}^{n+1}$. Была введена почти одновременно Г. Фубини [1] и Э. ІІтуди [2]. Ф.- Ш. м. задается формулой



\[

\left.d s^{2}=\frac{1}{|x|^{4}}\left(|x|^{2}|d x|^{2}-(x, \overline{d x}) \overline{(x}, d x\right)\right),

\]



где $|x|^{2}=(x, x) ;$ расстояние $\rho(\hat{x}, \hat{y})$ между точками $\widehat{x}=\mathbb{C} x, \hat{y}=\mathbb{C} y$, где $x, y \in \mathbb{C}^{n+1} \backslash\{0\}$, определяется из формулы



\[

\cos \rho \hat{x}, \hat{y})=\frac{|(x, y)|}{|x||y|}

\]



Ф.- Ш.м. является кәлеровой (и даже метрикой Ходжа), соответствующая ей кәлерова форма имеет вид



\[

\omega=\frac{i}{2 \pi} \partial \bar{\partial} \ln |z|^{2} .

\]



Ф.- ШІ. м.- это единственная с точностью до пропорциональности риманова метрика на $\mathbb{C} P^{n}$, инвариантная относительно унитарной группы $U(n+1)$, сохраняющей скалярное произведение. Пространство $\mathbb{C} P^{n}$, снабженное Ф.- ІІІ. м., является компактным әрмитовым симметрич. пространством ранга 1. Оно наз. также э л ли итическим әрмитовым пространством.



Jum.: [1] F u b in i G., "Atti Istit. Veneto», 1904, v. 63, p. 502-13; [2] S t u d y E., "Math. Ann.", 1905, Bd 60, S. $321-$ 78; [3] C a r t a n E., Leçons sur la géometrie projective complexe, Р., 1950; [4] X е лг а с о н С., Дифференциальная геометрия ч ж е н ь III ә н-ш ә н ь, Комшлексные многообразия, пер. с англ., М., 1961. A. Л. Онищик.



ФУКСА УРАВНЕНИЕ, у р а в н е н и е к л а с с а Ф у к с а - линейное однородное обыкновенное дифференциальное уравнение в комплексной области



\[

w^{(n)}+p_{1}(z) w^{(n-1)}+\ldots+p_{n}(z) w=0

\]



с аналитич. коэфффициентами, все особые точки к-рого на Римана сфере являются регулярными особыми точками. Для того чтобы уравнение (1) принадлежало классу Фукса, необходимо и достаточно, чтобы его коэфффициенты имели вид



\[

p_{j}(z)=\prod_{m=1}^{k}\left(z-z_{m}\right)^{-j} q_{j}(z)

\]



где $z_{1}, \ldots, z_{k}$-различные точки, $q_{j}(z)$-многочлен степени $\leqslant j(k-1)$. Система $w^{\prime}=A(z) w$ из $n$ уравнений принадлежит классу Фукса, если она имеет вид



\[

\frac{d w}{d z}=\sum_{m=1}^{k} \frac{A_{m}}{z-z_{m}} w

\]



где $z_{1}, \ldots, z_{k}$-различные точки, $A_{m} \neq 0$-постоянные матрицы порядка $n \times n$. Особыми для уравнения (1) и системы (2) являются точки $z_{1}, \ldots, z_{k}, \infty$. Для Ф. у. (1) справедливо то ждест во Ф у кс а:



\[

\sum_{j=1}^{n}\left(\sum_{m=1}^{k} \rho_{j}^{m}+\rho_{j}^{\infty}\right)=(k-1) \frac{n(n-1)}{2},

\]



где $\rho_{1}^{m}, \ldots, \rho_{n}^{m}$-характеристич. показатели в тотке $z_{m}$, а $\rho_{1}^{\infty}, \ldots, \rho_{n}^{\infty}-$ в точке $\infty$. Ф. у. (и системы) наз. также регуля рными уравнениями (системами). Этот класс уравнений и систем был введен Л. Фук$\operatorname{com}[1]$.



Пусть $D-$ сфера Римана с проколами в точках $z_{1}, \ldots$ $\ldots, z_{k}, \infty$. Любое нетривиальное решение Ф. у. (1) (соответственно любая компонента решения системы (2)) есть аналитическая в области $D$ функция. Как правило, эта функция бесконечнозначна, а все особые точки уравнения (1) (системы (2)) являются ее точками ветвления бесконечного порядка.



Ф. у. 2-го порядка с особыми точками $z_{1}, z_{2}, \ldots, z_{k}, \infty$ имеет вид



\[

\begin{gathered}

w^{\prime \prime}+\sum_{m=1}^{k} \frac{1-\left(\rho_{1}^{m}+\rho_{2}^{m}\right)}{z-z_{m}} w^{\prime}+ \\

+\sum_{m=1}^{k}\left[\frac{\rho_{1}^{m} \rho_{2}^{m} \prod_{j=1}^{\prime k}\left(z_{m}-z_{j}\right)}{z-z_{m}}+Q_{k-2}(z)\right] \times \\

\times \frac{w}{\prod_{m=1}^{k}\left(z-z_{m}\right)}=0

\end{gathered}

\]



где $Q_{k-2}(z)-$ многочлен степени $k-2$. Преобразование $w=\left(z-z_{m}\right)^{l} \tilde{w}$ переводит Ф. у. в Ф. у., причем



\[

\begin{aligned}

& \left(\rho_{1}^{m}, \rho_{2}^{m}\right) \longrightarrow\left(\rho_{1}^{m}-l, \rho_{2}^{m}-l\right), \\

& \left(\rho_{1}^{\infty}, \rho_{2}^{\infty}\right) \longrightarrow\left(\rho_{1}^{\infty}+l, \rho_{2}^{\infty}+l\right),

\end{aligned}

\]





\end{document}